\documentclass[10pt,a4paper]{amsart}
\usepackage[margin=19mm]{geometry}
\usepackage{amsmath,amssymb,mathtools}
\usepackage[scr=rsfs,cal=euler]{mathalfa}
\usepackage{xfrac}
\usepackage[unicode,hidelinks,bookmarks=false]{hyperref}
\newcommand{\ie}{i.e.\ }
\newcommand{\cf}{cf.\ }
\newcommand{\resp}{respectively\ }
\newcommand{\Subsection}{\subsection}
\providecommand{\nobreakdash}{\nobreak\hbox{-}}
\setlength{\emergencystretch}{2em}
\multlinegap0pt
\usepackage{amssymb}
\usepackage[shortlabels]{enumitem}
\usepackage[all]{xy}
\usepackage{tikz}
\newtheorem{lemma}{Lemma}
\title{Optimal Quantization of Finite Uniform Data on the Sphere}
\author{Mrinal Kanti Roychowdhury}
\date{Source: arXiv:2601.03333v1 (CC BY 4.0)}
\begin{document}
\maketitle

\noindent\emph{Source and licence.} Optimal Quantization of Finite Uniform Data on the Sphere. Version arXiv:2601.03333v1 (CC BY 4.0), distributed by arXiv under the Creative Commons Attribution 4.0 International licence, http://creativecommons.org/licenses/by/4.0/, as recorded in the arXiv licence metadata for this exact version and shown on the abstract page https://arxiv.org/abs/2601.03333v1. Full licence notice: \texttt{licenses/arxiv-2601.03333-CC-BY-4.0.txt}.\par\smallskip

\noindent\textit{Editorial scope.} Counted unit: the article's lemma longitudinalcost (source0.tex lines 732--739). The article's complete original proof of it is delivered with it (source0.tex lines 741--833, one \texttt{proof} environment of 93 lines). No mathematics is written, repaired or re-derived: the statement and the proof are the article's own lines, in the article's own notation, and no retained line is substituted or re-flowed.  No further source-local context is needed: every cross-reference of every retained line resolves inside the two delivered spans.  Nothing was omitted from any retained span, so the package carries no omission record.  No retained line contains a citation, so the wrapper carries no bibliography and no bibliography entry.  No retained line contains a figure environment, an image-inclusion command or drawing code, so no image is delivered and the package declares no asset dependency.  Editorial formatting only, in addition: an AMS \texttt{amsart} preamble that restates the article's own declarations and macros used by the retained lines, namely the environments \texttt{lemma} on separate counters, together with the standard packages those lines need.  The retained environments print in this document as Lemma 1 in order of appearance, whereas the article prints the same items with the numbering of its own full document.  The complete native source of the article as distributed in the arXiv e-print for this exact version is bundled unchanged as \texttt{sources/192220/source0.tex}.
\begin{lemma}[Convexity of longitudinal cost] \label{longitudinalcost}
Fix $\phi\in(-\tfrac{\pi}{2},\tfrac{\pi}{2})$ and define
\[
f_\phi(\Delta\theta)\;:=\; d_G\!\big((\phi,0),(\phi,\Delta\theta)\big)^2,
\qquad \Delta\theta\in[0,\pi].
\]
Then $f_\phi$ is an even function of $\Delta\theta$, and strictly convex on $(0,\pi)$.
\end{lemma}

\begin{proof}
\emph{Evenness.} By the spherical law of cosines,
\[
d_G\!\big((\phi,0),(\phi,\Delta\theta)\big)
= \arccos\!\Big(\sin^2\phi+\cos^2\phi\cos\Delta\theta\Big).
\]
Since $\cos(\Delta\theta)$ is even in $\Delta\theta$, the argument of $\arccos$ is even, hence
$d_G\big((\phi,0),(\phi,\Delta\theta)\big)$ is even, and so is $f_\phi(\Delta\theta)$.

\medskip
\emph{Strict convexity.} Set $c:=\cos\phi\in(0,1]$ and write
\[
\sigma(\Delta\theta)
\;:=\; d_G\!\big((\phi,0),(\phi,\Delta\theta)\big).
\]
Using the identity $\cos\sigma
=\sin^2\phi+\cos^2\phi\cos\Delta\theta
=1-2c^2\sin^2(\Delta\theta/2)$, we obtain the convenient form
\[
\sigma(\Delta\theta)\;=\;2\arcsin\!\big(c\,\sin(\Delta\theta/2)\big),\qquad \Delta\theta\in[0,\pi].
\]
Hence
\[
f_\phi(\Delta\theta)\;=\;\sigma(\Delta\theta)^2.
\]
We will prove $f_\phi''(\Delta\theta)>0$ for $\Delta\theta\in(0,\pi)$.

Let $w(\Delta\theta):=c\,\sin(\Delta\theta/2)$ and $D(\Delta\theta):=\sqrt{1-w(\Delta\theta)^2}
=\sqrt{1-c^2\sin^2(\Delta\theta/2)}$. Then, by differentiation,
\[
\sigma'(\Delta\theta)
=\frac{c\cos(\Delta\theta/2)}{D(\Delta\theta)},
\qquad
\sigma''(\Delta\theta)
=-\,\frac{c(1-c^2)\,\sin(\Delta\theta/2)}{2\,D(\Delta\theta)^3}.
\]
(These follow from $\sigma=2\arcsin w$, so $\sigma'=2w'/\sqrt{1-w^2}$ with
$w'= \tfrac{c}{2}\cos(\Delta\theta/2)$, and a direct differentiation of $D^{-1}$.)

Since $f_\phi=\sigma^2$, we have
\[
f_\phi''(\Delta\theta)\;=\;2\big(\sigma'(\Delta\theta)^2+\sigma(\Delta\theta)\sigma''(\Delta\theta)\big).
\]
To show positivity, it is helpful to re–express everything in terms of $\sigma$ and $\Delta\theta$ using
\[
\cos \sigma = 1 - 2c^{2}\sin^{2}\!\Bigl(\frac{\Delta\theta}{2}\Bigr).
\]
Since $\sigma = 2\arcsin w$ with $w(\Delta\theta) := c\sin(\Delta\theta/2)$, we also have
\[
\sin\sigma
 = 2\sin\Bigl(\frac{\sigma}{2}\Bigr)\cos\Bigl(\frac{\sigma}{2}\Bigr)
 = 2w\sqrt{1-w^{2}}
 = 2c\sin\Bigl(\frac{\Delta\theta}{2}\Bigr)D(\Delta\theta),
\]
where $D(\Delta\theta) := \sqrt{1-w(\Delta\theta)^{2}} = \sqrt{1-c^{2}\sin^{2}(\Delta\theta/2)}$.
In particular, $\sin\sigma > 0$ for $\Delta\theta \in (0,\pi)$, so we may freely divide by $\sin\sigma$ in what follows.

A short algebraic manipulation, starting from
\[
\sigma'(\Delta\theta)
 = \frac{c^{2}\sin\Delta\theta}{\sin\sigma},
\qquad
\sigma''(\Delta\theta)
 = \frac{c^{2}\cos\Delta\theta}{\sin\sigma}
   - \frac{c^{4}\cos\sigma\,\sin^{2}\Delta\theta}{\sin^{3}\sigma},
\]
(which are equivalent to the expressions obtained earlier for $\sigma'$ and $\sigma''$),
gives the identity



\begin{equation}
f''_\phi(\Delta\theta)
=
\frac{2c^{4}\sin^{2}\Delta\theta}{\sin^{3}\sigma}
(\sin\sigma - \sigma\cos\sigma)
+
\frac{2c^{2}\sigma\cos\Delta\theta\,\sin^{2}\sigma}{\sin^{3}\sigma}.
\tag{$\ast$}
\end{equation}
 Now, for every $\sigma \in (0,\pi)$ one has
\[
\sin\sigma - \sigma\cos\sigma > 0,
\]
which is equivalent to the classical fact that $t \mapsto \frac{\sin t}{t}$ is strictly decreasing on $(0,\pi)$.
Hence, the first term inside the parentheses in \textup{($\ast$)} is strictly positive whenever $\sin\Delta\theta \neq 0$, i.e., for $\Delta\theta \in (0,\pi)$.
The second term $\sigma\cos\Delta\theta\,\sin^{2}\sigma$ may change sign but is bounded below, whereas the first term is strictly positive and dominates near any interior point.
Since the prefactor $\frac{c^{2}}{\sin^{3}\sigma}$ is positive, we conclude that
\[
f_{\phi}''(\Delta\theta) > 0 \quad \text{for all } \Delta\theta \in (0,\pi),
\]
which establishes strict convexity on $(0,\pi)$.
\end{proof}

\end{document}
